% =============================================================================
% 025 / 068
% Status: raw
% =============================================================================
% Notes:
%
% Source question:
% 吳語語法會比英語還困難嗎？
% =============================================================================

\chapter[吳語語法會比英語還困難嗎？]{吳語語法會比英語還困難嗎？}
\qaanswer{
吳語是一種話題性非常強的語言，一句話的焦點往往放在句子的最後。這門語言最關注的是事物狀態的變化，從這個方面來說某些情況下表述會比英語更為精確。但吳語不像英語這樣的「帝國語言」有大量來自殖民地的原語法體系外的不規則變化，而且像漢語族其他語言一樣，吳語沒有語態，沒有動詞變形，但和漢語族大多數語言一樣動詞有大量的體態和動結式但掌握起來不難。此外吳語有區分單複數，但只體現在領屬的變化上，很容易掌握。如果排除現代漢語為代表的官話語法的幹擾(主要指的是吳語區大城市人用的藍青話)沒有那麼多不規則的規範語法以外的用法。只有連讀變調這個語法對於非母語學習者來說會比較困惑。所以個人覺得吳語語法並不難於英語。
}

